\documentclass[11pt]{article}
\usepackage[a4paper,margin=2cm,top=1.8cm,bottom=1.8cm]{geometry}
\usepackage{fontspec}
\setmainfont{Carlito}
\usepackage{xcolor}
\usepackage{booktabs}
\usepackage{amssymb}
\usepackage{tabularx}
\usepackage{array}
\usepackage{enumitem}
\usepackage{listings}
\usepackage[most]{tcolorbox}
\usepackage{fancyhdr}
\usepackage{lastpage}
\usepackage{hyperref}

\definecolor{azulOscuro}{RGB}{14,47,90}
\definecolor{azulMedio}{RGB}{31,92,153}
\definecolor{tealUteq}{RGB}{14,140,127}
\definecolor{grisFila}{RGB}{242,246,250}
\definecolor{codeBg}{RGB}{246,248,250}
\definecolor{codeBorde}{RGB}{200,210,222}
\definecolor{comentario}{RGB}{90,110,90}
\definecolor{cadena}{RGB}{160,60,30}
\definecolor{clave}{RGB}{20,70,140}

\hypersetup{colorlinks=true,linkcolor=azulMedio,urlcolor=tealUteq}

\setlength{\parindent}{0pt}
\setlength{\parskip}{4pt}
\renewcommand{\arraystretch}{1.2}

\pagestyle{fancy}
\fancyhf{}
\renewcommand{\headrulewidth}{0.6pt}
\renewcommand{\headrule}{\hbox to\headwidth{\color{azulMedio}\leaders\hrule height \headrulewidth\hfill}}
\fancyhead[L]{\footnotesize\color{azulMedio}\textbf{GA-SUM-03 / PE-U2} -- Guia de resolucion paso a paso}
\fancyhead[R]{\footnotesize\color{azulMedio}UTEQ -- Aplicaciones Distribuidas}
\fancyfoot[C]{\footnotesize\color{azulMedio}Pagina \thepage\ de \pageref{LastPage}}

\lstdefinestyle{base}{
  basicstyle=\ttfamily\scriptsize,
  backgroundcolor=\color{codeBg},
  frame=single,
  framesep=4pt,
  rulecolor=\color{codeBorde},
  breaklines=true,
  breakatwhitespace=true,
  columns=fullflexible,
  keepspaces=true,
  showstringspaces=false,
  tabsize=2,
  commentstyle=\color{comentario},
  stringstyle=\color{cadena},
  keywordstyle=\color{clave}\bfseries,
  aboveskip=6pt,
  belowskip=6pt
}
\lstset{style=base}

\lstdefinelanguage{shell}{
  morecomment=[l]{\#},
  morestring=[b]",
  morestring=[b]',
  sensitive=true
}
\lstdefinelanguage{yamlx}{
  morecomment=[l]{\#},
  morestring=[b]",
  morestring=[b]',
  sensitive=true
}

\newcommand{\pasohead}[2]{%
\vspace{4pt}%
\begin{tcolorbox}[colback=azulMedio,colframe=azulMedio,boxrule=0pt,arc=2pt,left=8pt,right=8pt,top=4pt,bottom=4pt]
\textcolor{white}{\textbf{Paso #1.\ \ #2}}\end{tcolorbox}}

\newtcolorbox{nota}{colback=tealUteq!8,colframe=tealUteq,boxrule=0.7pt,arc=2pt,left=6pt,right=6pt,top=4pt,bottom=4pt}

\begin{document}

% ===== PORTADA =====
\begin{center}
{\LARGE\textbf{\textcolor{azulOscuro}{UNIVERSIDAD TECNICA ESTATAL DE QUEVEDO}}}\\[2pt]
{\large\textcolor{tealUteq}{Facultad de Ciencias de la Computacion y Diseno Digital}}\\[2pt]
{\normalsize\textcolor{azulMedio}{Software (Redisenio) -- Aplicaciones Distribuidas [20701] -- 7mo Nivel A}}\\[6pt]
\begin{tcolorbox}[colback=azulOscuro,colframe=azulOscuro,boxrule=0pt,arc=2pt,left=10pt,right=10pt,top=6pt,bottom=6pt,width=\textwidth]
\centering\textcolor{white}{\Large\textbf{Guia de resolucion paso a paso}}\\[2pt]
\textcolor{white}{GA-SUM-03 / PE-U2 -- Arquitectura, API y Contenerizacion del Sistema}
\end{tcolorbox}
\end{center}

\textbf{\textcolor{azulOscuro}{Objetivo de esta guia.}} Construir, contenerizar y demostrar un sistema de microservicios REST con autenticacion JWT, API Gateway (Nginx) y Docker Compose, que constituye la Entrega 2 del PFC. La solucion de referencia usa \textbf{FastAPI (Python)}; los mismos pasos aplican a Spring Boot (Java) o Express (Node.js), que tambien autoriza la guia. Cada equipo debe adaptar el ``recurso principal'' a su propio PFC (producto, estudiante, laboratorio, ticket, etc.).

\begin{nota}
\small \textbf{Resultado esperado.} Al ejecutar \texttt{docker compose up -{}-build} deben quedar operativos: tres microservicios con base de datos propia, un API Gateway en el puerto 8080 y un servicio de monitoreo (Portainer) en el 9000, accesibles y verificables con Postman.
\end{nota}

% ===== ARQUITECTURA OBJETIVO =====
\section*{\textcolor{azulOscuro}{1.\ Arquitectura objetivo}}
El sistema enruta todas las peticiones del cliente a traves del API Gateway, que las distribuye a cada microservicio. Cada microservicio es independiente y posee su propia base de datos (patron \emph{Database per Service}).

\begin{lstlisting}
                         +-------------------+
        Cliente  ---->   |  API Gateway      |  (Nginx, puerto 8080)
        (Postman)        |  rate limit+log   |
                         +----+----+----+----+
              /api/auth/ |    /api/resources/ |  /api/notifications
                         v         v          v
            +-----------------+ +----------------+ +---------------------+
            |  auth-service   | | resource-svc   | | notification-service|
            |  :8001  JWT     | | :8002  CRUD    | | :8003  eventos      |
            +--------+--------+ +-------+--------+ +----------+----------+
                     |                  |                     |
                 users.db          resources.db         notifications.db
\end{lstlisting}

El diagrama C4 nivel 2 (contenedores) y nivel 3 (componentes) se elabora en draw.io o Structurizr y se incluye como figura en el documento LaTeX. El bosquejo anterior solo orienta; no sustituye al C4 formal exigido.

% ===== ESTRUCTURA DEL REPO =====
\section*{\textcolor{azulOscuro}{2.\ Estructura del repositorio}}
Crea el repositorio en GitHub \textbf{publico} y respeta esta estructura. La URL debe quedar escrita en el documento de la entrega y el ultimo commit debe corresponder a la fecha y hora de entrega.

\begin{lstlisting}
pfc-e2/
  docs/api/            # contratos OpenAPI 3.0 (YAML) exportados de Swagger
    auth.yaml
    resource.yaml
    notification.yaml
  auth-service/        # main.py, requirements.txt, Dockerfile
  resource-service/
  notification-service/
  api-gateway/         # nginx.conf
  docker-compose.yml
  .env.example         # variables SIN valores reales
  .gitignore           # debe incluir .env y *.db
  informe/             # documento LaTeX de la Entrega 2
  README.md
\end{lstlisting}

\pasohead{0}{Preparacion del entorno (PowerShell por defecto)}
Instala Docker Desktop 4.x, Python 3.11+ y Git. Verifica las versiones antes de empezar.
\begin{lstlisting}[language=shell]
# Windows PowerShell (por defecto en esta guia)
docker --version
docker compose version
python --version
git --version
# Linux/macOS: los mismos comandos funcionan en bash/zsh
\end{lstlisting}

% ===== PASO 1 =====
\pasohead{1}{Diseno de la arquitectura y contratos OpenAPI}
Define los tres microservicios y sus responsabilidades: \texttt{auth-service} emite y valida JWT; \texttt{resource-service} expone el CRUD del recurso principal del PFC y verifica el token en cada operacion; \texttt{notification-service} registra un evento cada vez que se crea o modifica un recurso.

Para cada microservicio escribe el contrato \textbf{OpenAPI 3.0} en \href{https://editor.swagger.io}{Swagger Editor} y exporta el YAML en \texttt{docs/api/}. A continuacion, un contrato minimo de \texttt{auth-service} como modelo.

\begin{lstlisting}[language=yamlx]
openapi: 3.0.3
info:
  title: auth-service
  version: "1.0.0"
paths:
  /auth/login:
    post:
      summary: Autentica al usuario y devuelve un JWT
      requestBody:
        required: true
        content:
          application/json:
            schema:
              type: object
              required: [username, password]
              properties:
                username: { type: string }
                password: { type: string, format: password }
      responses:
        "200":
          description: Login correcto (JWT en data.access_token)
        "401":
          description: Credenciales invalidas
  /auth/validate:
    get:
      summary: Verifica un token Bearer
      security: [ { bearerAuth: [] } ]
      responses:
        "200": { description: Token valido }
        "401": { description: Token invalido o expirado }
components:
  securitySchemes:
    bearerAuth: { type: http, scheme: bearer, bearerFormat: JWT }
\end{lstlisting}

% ===== PASO 2 =====
\pasohead{2}{Implementacion de los microservicios (FastAPI)}
Todos los microservicios responden con el sobre JSON \texttt{\{status, data, message, timestamp\}}. El archivo \texttt{requirements.txt} (identico para los tres, ajustando solo el puerto en el Dockerfile) usa versiones verificadas en PyPI.

\begin{lstlisting}
# requirements.txt
fastapi==0.115.0
uvicorn[standard]==0.30.6
sqlalchemy==2.0.34
python-jose[cryptography]==3.3.0
passlib[bcrypt]==1.7.4
requests==2.32.3
pydantic==2.9.2
\end{lstlisting}

\textbf{auth-service/main.py} -- registro, login (JWT con expiracion de 1 h) y validacion.
\begin{lstlisting}[language=Python]
import os, datetime as dt
from fastapi import FastAPI, HTTPException, Header
from pydantic import BaseModel
from sqlalchemy import create_engine, Column, Integer, String
from sqlalchemy.orm import declarative_base, sessionmaker
from passlib.context import CryptContext
from jose import jwt, JWTError

SECRET  = os.getenv("JWT_SECRET", "change-me")
ALGO    = "HS256"
EXP_MIN = int(os.getenv("JWT_EXP_MIN", "60"))
DB_URL  = os.getenv("DATABASE_URL", "sqlite:///./users.db")

engine  = create_engine(DB_URL, connect_args={"check_same_thread": False})
Session = sessionmaker(bind=engine)
Base    = declarative_base()
pwd     = CryptContext(schemes=["bcrypt"], deprecated="auto")

class User(Base):
    __tablename__ = "users"
    id = Column(Integer, primary_key=True)
    username = Column(String, unique=True, index=True)
    password_hash = Column(String)
Base.metadata.create_all(engine)

app = FastAPI(title="auth-service")

def envelope(data=None, message="OK", status="success"):
    return {"status": status, "data": data, "message": message,
            "timestamp": dt.datetime.utcnow().isoformat() + "Z"}

class Credentials(BaseModel):
    username: str
    password: str

@app.post("/auth/register")
def register(c: Credentials):
    db = Session()
    if db.query(User).filter_by(username=c.username).first():
        raise HTTPException(status_code=409, detail="Usuario ya existe")
    db.add(User(username=c.username, password_hash=pwd.hash(c.password)))
    db.commit()
    return envelope({"username": c.username}, "Usuario registrado")

@app.post("/auth/login")
def login(c: Credentials):
    db = Session()
    u = db.query(User).filter_by(username=c.username).first()
    if not u or not pwd.verify(c.password, u.password_hash):
        raise HTTPException(status_code=401, detail="Credenciales invalidas")
    exp = dt.datetime.utcnow() + dt.timedelta(minutes=EXP_MIN)
    token = jwt.encode({"sub": c.username, "exp": exp}, SECRET, algorithm=ALGO)
    return envelope({"access_token": token, "token_type": "bearer",
                     "expires_in": EXP_MIN * 60}, "Login correcto")

@app.get("/auth/validate")
def validate(authorization: str = Header(default="")):
    token = authorization.replace("Bearer ", "")
    try:
        payload = jwt.decode(token, SECRET, algorithms=[ALGO])
    except JWTError:
        raise HTTPException(status_code=401, detail="Token invalido o expirado")
    return envelope({"username": payload["sub"]}, "Token valido")
\end{lstlisting}

\textbf{resource-service/main.py} -- CRUD con verificacion de JWT en cada endpoint; al crear o modificar avisa a \texttt{notification-service}.
\begin{lstlisting}[language=Python]
import os, datetime as dt, requests
from fastapi import FastAPI, HTTPException, Header, Depends
from pydantic import BaseModel
from sqlalchemy import create_engine, Column, Integer, String
from sqlalchemy.orm import declarative_base, sessionmaker
from jose import jwt, JWTError

SECRET    = os.getenv("JWT_SECRET", "change-me")
ALGO      = "HS256"
DB_URL    = os.getenv("DATABASE_URL", "sqlite:///./resources.db")
NOTIF_URL = os.getenv("NOTIFICATION_URL", "http://notification-service:8003")

engine  = create_engine(DB_URL, connect_args={"check_same_thread": False})
Session = sessionmaker(bind=engine)
Base    = declarative_base()

class Resource(Base):
    __tablename__ = "resources"
    id = Column(Integer, primary_key=True)
    name = Column(String)
    description = Column(String)
Base.metadata.create_all(engine)

app = FastAPI(title="resource-service")

def envelope(data=None, message="OK", status="success"):
    return {"status": status, "data": data, "message": message,
            "timestamp": dt.datetime.utcnow().isoformat() + "Z"}

def require_user(authorization: str = Header(default="")):
    token = authorization.replace("Bearer ", "")
    try:
        return jwt.decode(token, SECRET, algorithms=[ALGO])["sub"]
    except JWTError:
        raise HTTPException(status_code=401, detail="Token requerido o invalido")

class ResourceIn(BaseModel):
    name: str
    description: str = ""

def notify(event, rid, user):
    # Degradacion controlada: si notification-service no responde, no se cae el CRUD
    try:
        requests.post(f"{NOTIF_URL}/notifications",
                      json={"event": event, "resource_id": rid, "user": user},
                      timeout=2)
    except requests.RequestException:
        pass

@app.get("/resources")
def list_resources(user: str = Depends(require_user)):
    db = Session()
    rows = db.query(Resource).all()
    return envelope([{"id": r.id, "name": r.name,
                      "description": r.description} for r in rows])

@app.post("/resources")
def create_resource(body: ResourceIn, user: str = Depends(require_user)):
    db = Session()
    r = Resource(name=body.name, description=body.description)
    db.add(r); db.commit(); db.refresh(r)
    notify("created", r.id, user)
    return envelope({"id": r.id}, "Recurso creado")

@app.put("/resources/{rid}")
def update_resource(rid: int, body: ResourceIn, user: str = Depends(require_user)):
    db = Session(); r = db.query(Resource).get(rid)
    if not r: raise HTTPException(status_code=404, detail="No encontrado")
    r.name, r.description = body.name, body.description; db.commit()
    notify("updated", rid, user)
    return envelope({"id": rid}, "Recurso actualizado")

@app.delete("/resources/{rid}")
def delete_resource(rid: int, user: str = Depends(require_user)):
    db = Session(); r = db.query(Resource).get(rid)
    if not r: raise HTTPException(status_code=404, detail="No encontrado")
    db.delete(r); db.commit()
    return envelope({"id": rid}, "Recurso eliminado")
\end{lstlisting}

\textbf{notification-service/main.py} -- registra cada evento recibido.
\begin{lstlisting}[language=Python]
import os, datetime as dt
from fastapi import FastAPI
from pydantic import BaseModel
from sqlalchemy import create_engine, Column, Integer, String, DateTime
from sqlalchemy.orm import declarative_base, sessionmaker

DB_URL  = os.getenv("DATABASE_URL", "sqlite:///./notifications.db")
engine  = create_engine(DB_URL, connect_args={"check_same_thread": False})
Session = sessionmaker(bind=engine)
Base    = declarative_base()

class Notification(Base):
    __tablename__ = "notifications"
    id = Column(Integer, primary_key=True)
    event = Column(String); resource_id = Column(Integer)
    user = Column(String); created_at = Column(DateTime, default=dt.datetime.utcnow)
Base.metadata.create_all(engine)

app = FastAPI(title="notification-service")

def envelope(data=None, message="OK", status="success"):
    return {"status": status, "data": data, "message": message,
            "timestamp": dt.datetime.utcnow().isoformat() + "Z"}

class NotificationIn(BaseModel):
    event: str; resource_id: int; user: str

@app.post("/notifications")
def create_notification(n: NotificationIn):
    db = Session()
    row = Notification(event=n.event, resource_id=n.resource_id, user=n.user)
    db.add(row); db.commit(); db.refresh(row)
    return envelope({"id": row.id}, "Notificacion registrada")

@app.get("/notifications")
def list_notifications():
    db = Session(); rows = db.query(Notification).all()
    return envelope([{"id": r.id, "event": r.event,
                      "resource_id": r.resource_id, "user": r.user} for r in rows])
\end{lstlisting}

% ===== PASO 3 =====
\pasohead{3}{API Gateway con Nginx}
El Gateway escucha en el puerto 8080, aplica \emph{rate limiting} de 100 peticiones por minuto e IP, registra todas las peticiones y enruta a cada microservicio. Guarda esto como \texttt{api-gateway/nginx.conf}.

\begin{lstlisting}
worker_processes auto;
events { worker_connections 1024; }

http {
    # Logging: fecha, metodo, ruta, IP, codigo de respuesta
    log_format apigw '$time_iso8601 $request_method $request_uri $remote_addr $status';
    access_log /var/log/nginx/access.log apigw;

    # Rate limiting: 100 peticiones por minuto por IP
    limit_req_zone $binary_remote_addr zone=apilimit:10m rate=100r/m;

    upstream auth          { server auth-service:8001; }
    upstream resources     { server resource-service:8002; }
    upstream notifications { server notification-service:8003; }

    server {
        listen 8080;

        location /api/auth/ {
            limit_req zone=apilimit burst=20 nodelay;
            proxy_pass http://auth/auth/;            # /api/auth/login -> /auth/login
        }

        # Coleccion (listar/crear): coincidencia exacta
        location = /api/resources {
            limit_req zone=apilimit burst=20 nodelay;
            proxy_pass http://resources/resources;
        }
        # Elemento individual: /api/resources/5 -> /resources/5
        location /api/resources/ {
            limit_req zone=apilimit burst=20 nodelay;
            proxy_pass http://resources/resources/;
        }

        location = /api/notifications {
            limit_req zone=apilimit burst=20 nodelay;
            proxy_pass http://notifications/notifications;
        }
    }
}
\end{lstlisting}

% ===== PASO 4 =====
\pasohead{4}{Contenerizacion con Docker Compose}
Cada microservicio usa un \textbf{Dockerfile multi-stage} (etapa \emph{builder} para instalar dependencias y etapa \emph{runtime} ligera). El siguiente es el de \texttt{auth-service}; para los otros dos solo cambia el \texttt{EXPOSE}/\texttt{-{}-port} a 8002 y 8003.

\begin{lstlisting}
# auth-service/Dockerfile
# ===== Stage 1: builder =====
FROM python:3.11-slim AS builder
WORKDIR /app
COPY requirements.txt .
RUN pip install --user --no-cache-dir -r requirements.txt

# ===== Stage 2: runtime =====
FROM python:3.11-slim
WORKDIR /app
COPY --from=builder /root/.local /root/.local
COPY . .
ENV PATH=/root/.local/bin:$PATH
EXPOSE 8001
CMD ["uvicorn", "main:app", "--host", "0.0.0.0", "--port", "8001"]
\end{lstlisting}

El \texttt{docker-compose.yml} levanta los tres microservicios, sus bases de datos (volumenes), el Gateway, el servicio de monitoreo \textbf{Portainer} y una red \texttt{bridge} interna.

\begin{lstlisting}[language=yamlx]
services:
  auth-service:
    build: ./auth-service
    environment:
      - JWT_SECRET=${JWT_SECRET}
      - JWT_EXP_MIN=${JWT_EXP_MIN}
      - DATABASE_URL=sqlite:////data/users.db
    volumes: [auth-data:/data]
    networks: [backend]

  resource-service:
    build: ./resource-service
    environment:
      - JWT_SECRET=${JWT_SECRET}
      - DATABASE_URL=sqlite:////data/resources.db
      - NOTIFICATION_URL=http://notification-service:8003
    depends_on: [notification-service]
    volumes: [resource-data:/data]
    networks: [backend]

  notification-service:
    build: ./notification-service
    environment:
      - DATABASE_URL=sqlite:////data/notifications.db
    volumes: [notif-data:/data]
    networks: [backend]

  api-gateway:
    image: nginx:alpine
    ports: ["8080:8080"]
    volumes:
      - ./api-gateway/nginx.conf:/etc/nginx/nginx.conf:ro
      - gateway-logs:/var/log/nginx
    depends_on: [auth-service, resource-service, notification-service]
    networks: [backend]

  portainer:                       # servicio de monitoreo
    image: portainer/portainer-ce:latest
    ports: ["9000:9000"]
    volumes:
      - /var/run/docker.sock:/var/run/docker.sock
      - portainer-data:/data
    networks: [backend]

networks:
  backend:
    driver: bridge

volumes:
  auth-data:
  resource-data:
  notif-data:
  gateway-logs:
  portainer-data:
\end{lstlisting}

Crea \texttt{.env.example} (se versiona, sin secretos) y un \texttt{.env} local real que \textbf{no} se sube (queda en \texttt{.gitignore}).
\begin{lstlisting}
# .env.example  (committear este; NO el .env real)
JWT_SECRET=pon-aqui-un-secreto-largo-y-aleatorio
JWT_EXP_MIN=60
\end{lstlisting}

Levanta y verifica el sistema:
\begin{lstlisting}[language=shell]
docker compose up --build          # construye y arranca todo
docker compose ps                  # los 5 servicios deben quedar "running"
# Portainer: abrir http://localhost:9000 para monitorear contenedores
\end{lstlisting}

% ===== PASO 5 =====
\pasohead{5}{Pruebas con Postman (3 flujos end-to-end)}
Ejecuta los tres flujos obligatorios contra el Gateway (\texttt{http://localhost:8080}). En PowerShell puedes usar \texttt{curl.exe}; en Postman crea una coleccion con una variable \texttt{token}.

\begin{lstlisting}[language=shell]
# Flujo 1: registro -> login -> capturar JWT
curl.exe -X POST "http://localhost:8080/api/auth/register" -H "Content-Type: application/json" -d "{\"username\":\"demo\",\"password\":\"Demo123*\"}"
curl.exe -X POST "http://localhost:8080/api/auth/login"    -H "Content-Type: application/json" -d "{\"username\":\"demo\",\"password\":\"Demo123*\"}"
# Copiar data.access_token de la respuesta y guardarlo como $TOKEN

# Flujo 2: sin token (401) -> con token (200) -> POST -> PUT -> DELETE
curl.exe "http://localhost:8080/api/resources"                                   # -> 401
curl.exe "http://localhost:8080/api/resources" -H "Authorization: Bearer $TOKEN" # -> 200
curl.exe -X POST "http://localhost:8080/api/resources" -H "Authorization: Bearer $TOKEN" -H "Content-Type: application/json" -d "{\"name\":\"Recurso 1\",\"description\":\"demo\"}"
curl.exe -X PUT    "http://localhost:8080/api/resources/1" -H "Authorization: Bearer $TOKEN" -H "Content-Type: application/json" -d "{\"name\":\"Recurso 1b\",\"description\":\"editado\"}"
curl.exe -X DELETE "http://localhost:8080/api/resources/1" -H "Authorization: Bearer $TOKEN"

# Flujo 3: crear recurso -> verificar la notificacion generada
curl.exe -X POST "http://localhost:8080/api/resources" -H "Authorization: Bearer $TOKEN" -H "Content-Type: application/json" -d "{\"name\":\"Recurso 2\",\"description\":\"con notif\"}"
curl.exe "http://localhost:8080/api/notifications"                                # debe listar el evento "created"
# Linux/macOS: reemplazar curl.exe por curl y usar comillas simples para el -d
\end{lstlisting}

Documenta los resultados con capturas y una tabla \texttt{booktabs} en el informe. Modelo (rellenar con tus mediciones reales):

\begin{center}\small
\begin{tabular}{lcccc}
\toprule
\textbf{Endpoint} & \textbf{Metodo} & \textbf{Cod. esperado} & \textbf{Cod. obtenido} & \textbf{Latencia (ms)} \\
\midrule
/api/auth/login        & POST   & 200 & 200 & 41 \\
/api/resources (s/token) & GET  & 401 & 401 & 8  \\
/api/resources         & POST   & 200 & 200 & 23 \\
/api/resources/1       & PUT    & 200 & 200 & 19 \\
/api/resources/1       & DELETE & 200 & 200 & 15 \\
/api/notifications     & GET    & 200 & 200 & 11 \\
\bottomrule
\end{tabular}
\end{center}

% ===== PASO 6 =====
\pasohead{6}{Documento LaTeX de la Entrega 2}
El informe (minimo 12 paginas, compila sin errores) debe incluir: fundamento teorico completo de la Unidad 3 (cuatro subtemas con sus referencias IEEE \textbf{verificadas con DOI}, tomadas de la bibliografia de la guia), el diagrama C4 como figura con \texttt{caption} y \texttt{label}, los contratos OpenAPI en \texttt{lstlisting} (YAML), los Dockerfiles comentados, el \texttt{nginx.conf}, y la tabla \texttt{booktabs} de resultados de pruebas. Incluye ademas la \textbf{declaracion de uso de IA} (que herramienta, como y para que) y la \textbf{URL del repositorio}.

% ===== PASO 7 =====
\pasohead{7}{Repositorio y commit de entrega}
Sube todo a GitHub y asegura que el \textbf{ultimo commit coincida con la fecha y hora de entrega}; este punto es eliminatorio en la rubrica.
\begin{lstlisting}[language=shell]
git init
git add -A
git commit -m "PFC E2: microservicios + gateway + docker-compose"
git branch -M main
git remote add origin https://github.com/<usuario>/pfc-e2.git
git push -u origin main
git log -1 --format=%cd        # la fecha mostrada debe ser la de entrega
\end{lstlisting}

% ===== RESPUESTAS DE ANALISIS =====
\section*{\textcolor{azulOscuro}{Respuestas modelo a las preguntas de analisis (Anexo)}}
\textbf{1. ?Por que cada microservicio debe tener su propia base de datos?} El patron \emph{Database per Service} garantiza el bajo acoplamiento: cada servicio controla su esquema, puede desplegarse y escalar de forma independiente y elige la tecnologia de persistencia que mejor se ajuste. Una base compartida reintroduce acoplamiento fuerte (el ``monolito distribuido''): un cambio de esquema afecta a varios servicios y se pierde la autonomia. El costo es que las consultas que cruzan servicios requieren composicion via API o eventos, y la consistencia pasa a ser \emph{eventual} en lugar de transaccional.

\textbf{2. ?Que ocurre si \texttt{resource-service} cae? Circuit Breaker.} Las peticiones quedarian colgadas o devolverian errores en cascada. El patron \emph{Circuit Breaker} envuelve las llamadas a ese servicio con tres estados: \emph{Closed} (pasa el trafico), \emph{Open} (tras N fallos consecutivos corta de inmediato y responde con un \emph{fallback} o 503 sin esperar el timeout) y \emph{Half-Open} (deja pasar una prueba para ver si se recupero). En FastAPI puede implementarse con una libreria como \texttt{pybreaker} alrededor del cliente HTTP, combinada con \emph{timeout} y reintentos acotados. En esta solucion, la funcion \texttt{notify()} ya aplica una degradacion controlada (timeout de 2 s y \texttt{pass} ante fallo) para que la caida de notificaciones no tumbe el CRUD.

\textbf{3. ?El API Gateway introduce un cuello de botella?} Si: agrega un salto de red (latencia adicional) y es un posible punto unico de fallo. Se mitiga escalando Nginx horizontalmente detras de un balanceador, habilitando \emph{keep-alive} hacia los upstreams, cacheo de respuestas idempotentes y dimensionando \texttt{worker\_processes} y \texttt{worker\_connections}. Midiendo con y sin Gateway se comprueba que el sobrecosto suele ser de pocos milisegundos, ampliamente compensado por la centralizacion de autenticacion, \emph{rate limiting}, enrutamiento y observabilidad.

\textbf{4. ?Como escalar \texttt{resource-service} a 3 replicas?} El servicio debe ser \emph{stateless} (el estado vive en su base de datos). Se escala con Compose y se deja que Nginx balancee por el nombre de servicio (DNS round-robin de Docker):
\begin{lstlisting}[language=shell]
docker compose up --build --scale resource-service=3
\end{lstlisting}
Requisitos: no fijar \texttt{container\_name} ni publicar un puerto de host en \texttt{resource-service} (para evitar colisiones), y en Nginx mantener el \texttt{upstream resources} apuntando al nombre del servicio; Docker resolvera las 3 instancias. Si se requiere control fino, se usa \texttt{deploy.replicas} con Swarm o se migra a Kubernetes (Deployment con 3 \emph{pods} y un \emph{Service}).

% ===== CHECKLIST =====
\section*{\textcolor{azulOscuro}{Lista de verificacion final (mapeada a la rubrica)}}
\begin{center}\small
\begin{tabularx}{\textwidth}{cX c}
\toprule
\textbf{Dim.} & \textbf{Verifica que...} & \textbf{Hecho} \\
\midrule
D1 & Diagrama C4 nivel 2 y 3; contratos OpenAPI 3.0 (YAML) en \texttt{docs/api/} & $\square$ \\
D2 & Los 3 microservicios responden; JWT emitido y verificado (Bearer); Gateway con rate limit y logging & $\square$ \\
D3 & Dockerfile multi-stage; \texttt{docker compose up} sin errores; Portainer operativo; \texttt{.env.example} sin secretos & $\square$ \\
D4 & Los 3 flujos Postman documentados con capturas y tabla de latencias & $\square$ \\
D5 & Fundamento teorico U3 (4 subtemas, refs IEEE con DOI) y 4 preguntas respondidas & $\square$ \\
D6 & Informe LaTeX $\geq$12 pag.; URL de GitHub en el documento; ultimo commit = entrega; declaracion de IA & $\square$ \\
\bottomrule
\end{tabularx}
\end{center}

\begin{nota}
\small \textbf{Recordatorio de integridad.} Las referencias deben ser reales y verificables (DOI/ISBN que corresponda al titulo); las citas inventadas se penalizan drasticamente. Declara con honestidad el uso de IA: que herramienta, como y para que la usaste.
\end{nota}

\end{document}
